\section{Uniform marginal expansions}\label{sec:marginal-proof}

Throughout the proof, constants denoted \(C_q\) may change from line to
line but depend only on \(q\). The estimates below also explain why
the error is uniform when \(d=a-b\) grows.

\subsection{Expansion of the embedding product}

For an alternating \(j\times j\) form \(H\) of rank \(2s\), put
\[
 W(H)=\frac{q^{2(j-s)}-1}{q^2-1},
 \qquad Q_j=\frac{q^j-1}{q-1}.
\]

\begin{lemma}[Uniform product estimates]\label{lem:products}
Including the zero convention for infeasible embeddings,
\begin{align}
 0\le\Phi_D(j,s)&\le1,\qquad
 |\Phi_D(j,s)-1|\le q^{-D}W(H),\label{eq:product-first}\\
 \Phi_D(j,s)
   &=1-q^{-D}W(H)+O_q(q^{-2D}q^{4j}).\label{eq:product-second}
\end{align}
For \(j\le b\le a\), \(\Phi_{a+b}(j,s)\ge\eta_\infty\).
Moreover,
\begin{equation}\label{eq:Wmean}
 \E_HW(H)=Q_j.
\end{equation}
\end{lemma}

\begin{proof}
On the feasible range, the sum of the deficits of the factors in
\(\Phi_D\) is exactly
\[
 \sum_{r=D/2-j+s+1}^{D/2}q^{-2r}=q^{-D}W(H).
\]
For nonnegative \(t_i\le1\), their product satisfies
\[
 0\le \prod_i(1-t_i)-1+\sum_i t_i
       \le\frac12\left(\sum_i t_i\right)^2.
\]
For example, interpret the \(t_i\) as the probabilities of independent
events and apply the first two inclusion--exclusion bounds.
This proves the two estimates on the feasible range.
If \(j-s>D/2\), then \(q^{-D}W(H)\ge1\);
the estimates follow from \(\Phi_D=0\) and the quadratic bound
\(t-1\le t^2\) for \(t\ge1\).
The lower bound for \(\Phi_{a+b}\) follows because it is a finite
product of a subset of the factors of \(\eta_\infty\).

For \(j\ge1\), a fixed nonzero vector is in the radical of a uniform
alternating \(H\) with probability \(q^{-(j-1)}\).
Thus
\[
 \E_H q^{\nul H}=1+(q^j-1)q^{1-j}.
\]
Since \(2(j-s)=j+\nul H\), inserting this expression into the
definition of \(W\) proves \eqref{eq:Wmean}. The case \(j=0\) is
immediate.
\end{proof}

Use \(\Phi_b^2-\Phi_b=(\Phi_b-1)+(\Phi_b-1)^2\) rather than
multiplying two potentially large error estimates. Together with
\(\Phi_{a+b}\ge\eta_\infty\), this gives
\[
 \frac{\Phi_b^2}{\Phi_{a+b}}-\Phi_b
   =-vW+O_q(v^2q^{4j}+uvq^{2j}).
\]
Likewise,
\[
 \Phi_b(\Phi_a-1)=-uW+O_q(uvq^{4j}),
\]
because \(u\le v\). By \eqref{eq:unified},
\begin{align}
 \mu_j-\mu_j^0
   &=-v m_{j,d}Q_j+O_q(v^2m_{j,d}q^{4j}),
       \label{eq:rank-moment-error}\\
 \nu_j-\nu_j^0
   &=-u m'_jQ_j+O_q(uvm'_jq^{4j}).
       \label{eq:defect-moment-error}
\end{align}
At \(j=0\) both errors are exactly zero.

\subsection{Identification of the signed correction}

The identity
\[
 (q^k-q^j)\qbin k j
    =q^j(q^{j+1}-1)\qbin{k}{j+1}
\]
and the recurrences
\[
 (q^{j+1}-1)m_{j+1,d}=q^{-j-d}m_{j,d},
 \qquad
 (q^{j+1}-1)m'_{j+1}=q\,m'_j
\]
show from \eqref{eq:limit-moments} that
\[
 \sum_k q^k\qbin k j\pi_d(k)=(q^j+q^{-d})m_{j,d},
 \qquad
 \sum_k q^k\qbin k j\lambda(k)=(q+1)q^jm'_j.
\]
Consequently
\begin{equation}\label{eq:signed-moments}
 \sum_k\qbin k j\xi_d(k)=-m_{j,d}Q_j,
 \qquad
 \sum_k\qbin k j\zeta(k)=-m'_jQ_j.
\end{equation}
At \(j=0\) this proves zero total mass. The signs in
\eqref{eq:signed} yield
\begin{equation}\label{eq:signed-norms}
 \frac12\normone{\xi_d}
    =\xi_d(0)=\frac{q^{-d}\gamma_\infty}{(q-1)\gamma_d},
 \qquad
 \frac12\normone{\zeta}
    =\zeta(0)=\frac{q\theta_\infty}{q^2-1}.
\end{equation}

\subsection{An explicit inversion estimate}

Controlling finitely many moments is not by itself a bound on the
whole law. We use interpolation with estimates strong enough to
control every point mass and the infinite correction tails.

\begin{lemma}[Lagrange interpolation on geometric grids]\label{lem:interpolation}
For \(0\le k\le n\), put
\[
 B_{k,n}(x)=\prod_{\substack{0\le j\le n\\j\ne k}}
                   \frac{x-q^j}{q^k-q^j}.
\]
For every coefficient index \(i\),
\begin{equation}\label{eq:Lcoeff}
 |[x^i]B_{k,n}(x)|
    \le\gamma_\infty^{-3}
       q^{-k(k+1)/2}q^{-i(i-1)/2}.
\end{equation}
The polynomial \(B^{\mathrm{ev}}_{\ell,h}\) for the grid
\(1,q^2,\ldots,q^{2h}\) has
\begin{equation}\label{eq:Leven-coeff}
 |[x^i]B^{\mathrm{ev}}_{\ell,h}(x)|
    \le\eta_\infty^{-3}
       q^{-\ell(\ell+1)}q^{-i(i-1)}.
\end{equation}
For \(z>n\) and even \(z>2h\), respectively,
\begin{align}
 |B_{k,n}(q^z)|
    &\le\gamma_\infty^{-2}
       q^{nz-n(n+1)/2-k(k-1)/2},\label{eq:Ltail}\\
 |B^{\mathrm{ev}}_{\ell,h}(q^z)|
    &\le\eta_\infty^{-2}
       q^{hz-h(h+1)-\ell(\ell-1)}.\label{eq:Leven-tail}
\end{align}
\end{lemma}

\begin{proof}
Normalize each numerator factor as \(1-xq^{-j}\).
The absolute value of the constant prefactor is
\(q^{-k(k+1)/2}/(\gamma_k\gamma_{n-k})\).
The sum of the products of \(i\) distinct elements of
\(1,q^{-1},q^{-2},\ldots\) is
\(q^{-i(i-1)/2}/\gamma_i\):
remove the staircase \(0,1,\ldots,i-1\) from the increasing
indices and sum their nonnegative gaps.
Omitting one index can only decrease this nonnegative sum,
proving \eqref{eq:Lcoeff}. Replacing \(q\) by \(q^2\) proves
\eqref{eq:Leven-coeff}.
For \(z>n\), bound every numerator factor by \(q^z\) and use
\[
 \prod_{j\ne k}|q^k-q^j|
 =q^{n(n+1)/2+k(k-1)/2}\gamma_k\gamma_{n-k}.
\]
This gives \eqref{eq:Ltail}, and the even version follows in the same way.
\end{proof}

Counting linear maps from \(\F^i\) to \(\F^z\) by their image gives
\begin{equation}\label{eq:power-conversion}
 q^{iz}=\sum_{j=0}^i
          \left(\prod_{r=0}^{j-1}(q^i-q^r)\right)\qbin z j.
\end{equation}
This will turn Gaussian moments into the ordinary power moments
needed for interpolation.

\subsection{Proof of the rank expansion}

Set \(\rho=\Law(K_A)-\Law(K_U)-v\xi_d\).
By \eqref{eq:rank-moment-error} and \eqref{eq:signed-moments}, its \(j\)th Gaussian
moment vanishes at \(j=0\), and for \(1\le j\le b\) has absolute value
at most
\[
 C_q v^2q^{-j^2-dj+4j}
 \le C_q uvq^{-j^2+4j}.
\]
Here is the essential uniformity step:
\(q^{-dj}\le q^{-d}\) when \(j\ge1\), and \(v^2q^{-d}=uv\).
For every \(i\le b\), \eqref{eq:power-conversion} yields
\begin{equation}\label{eq:rank-power-error}
 \left|\int q^{iz}\,d\rho(z)\right|
   \le C_q uv\sum_{j=1}^i q^{-j^2+(i+4)j}
   \le C'_q uvq^{(i+4)^2/4}.
\end{equation}
Pairing with \eqref{eq:Lcoeff} gives
\[
 \left|\int B_{k,b}(q^z)\,d\rho(z)\right|
 \le C_q uvq^{-k(k+1)/2},
\]
because
\(\sum_{i\ge0}q^{-i^2/4+5i/2+4}<\infty\).

Only the correction \(v\xi_d\) extends beyond the interpolation grid.
Uniformly in \(d\ge0\), for \(z\ge1\),
\[
 |\xi_d(z)|\le C_q q^{-dz-z^2+z}
             \le C_q q^{-d}q^{-z^2+z}.
\]
Combining this with \eqref{eq:Ltail}, setting \(z=b+1+t\), and summing
first over \(t\ge0\) and then over \(k\le b\), gives
\begin{equation}\label{eq:rank-tail}
 \sum_{k=0}^b\left|
    \sum_{z>b}B_{k,b}(q^z)\xi_d(z)\right|
 \le C_q q^{-d}q^{-b(b+1)/2}.
\end{equation}
Indeed, the remaining \(z\)-exponent is
\(-z^2+(b+1)z=-(b+1)t-t^2\).
Also
\[
 \sum_{z>b}|\xi_d(z)|
       \le C_q q^{-d}q^{-b(b+1)}.
\]
After multiplication by \(v\), both tails are \(O_q(uv)\).
Summing the interpolated estimates over \(k\) proves
\eqref{eq:rank-expansion}.

\subsection{Proof of the defect expansion}

Set
\(\rho'=\Law(L_U)-\Law(\nul T_b)-u\zeta\).
Its Gaussian-moment errors are bounded by
\(C_q uvm'_jq^{4j}\), with the zeroth error equal to zero.
For \(i\le b\), \eqref{eq:power-conversion} gives
\begin{equation}\label{eq:defect-power-error}
 \left|\int q^{iz}\,d\rho'(z)\right|
 \le C_q uvq^{i(i+9)/2}.
\end{equation}
For the sum giving this bound, the exponent
\(-j(j-1)/2+(i+4)j\) increases by at least five at each step
\(j\le i\), so the last term bounds the sum up to a constant
depending only on \(q\).

Put \(h=b/2\). Both finite laws and \(\zeta\) are supported on even
nullities, so \eqref{eq:Leven-coeff} applies. Pairing it with
\eqref{eq:defect-power-error} leaves the summable exponent
\(-i^2/2+11i/2\); hence
\[
 \left|\int B^{\mathrm{ev}}_{\ell,h}(q^z)\,d\rho'(z)\right|
    \le C_q uvq^{-\ell(\ell+1)}.
\]
For even \(z\ge2\),
\[
 |\zeta(z)|\le C_q q^{-z(z-3)/2}.
\]
Using \eqref{eq:Leven-tail} and setting \(z=2h+2+2t\) gives
\begin{equation}\label{eq:even-tail}
 \sum_{\ell=0}^h\left|
   \sum_{\substack{z>b\\z\ {\rm even}}}
        B^{\mathrm{ev}}_{\ell,h}(q^z)\zeta(z)\right|
     \le C_q q^{1-h^2}\le C'_q q^{-b}.
\end{equation}
The exponent at \(t=0\) is \(1-h^2\) and subsequently decreases
quadratically. The final inequality holds for \(h\ge3\);
the two cases \(h=1,2\) are absorbed by a constant depending only on \(q\).
The mass of \(\zeta\) beyond the grid is bounded by
\[
 C_q q^{-(b+2)(b-1)/2}\le C'_q q^{-b}.
\]
Multiplying the tails by \(u\), then summing the interpolation estimates,
proves \eqref{eq:defect-expansion}.

Finally, if \(P-Q=t\sigma+R\), then
\[
 \left|\TV(P,Q)-\frac{|t|}{2}\normone{\sigma}\right|
       \le\frac12\normone R.
\]
Apply this inequality and \eqref{eq:signed-norms} to obtain
\eqref{eq:rank-tv} and \eqref{eq:defect-tv}, completing
\cref{thm:marginals}.

